% GPT-6 Astra Pro; effort not exposed; /root/replacement; 22 November 2026.
\section{HD-F1-R15: a line-character endpoint cannot be coupled to a rank-two or rank-three bundle with balanced difference}
\label{sec:followup-rm-low-rank-adjacent}

There are no actual vector bundle $A$ of rank $r=2$ or $3$
and coherent rank-one sheaf $T$ with
$\operatorname{ch}(T)=e^{\ell}$ satisfying
\[
 \operatorname{ch}(A)-\operatorname{ch}(T)
        =(r-1)e^{-B}(1-qA_1)(1-qA_2),\qquad q>0.
 \tag{F1R15.1}
\]
The slope $\ell$ and $B$ may be arbitrary rational algebraic
divisor classes. No semihomogeneity of $A$, finite action
or semiregularity is imposed. In particular every actual
adjacent two-cohomology object with these endpoints fails
the required whole-character equation, regardless of its
connecting morphism in $\operatorname{Ext}^2(T,A)$.
Every other actual two-cohomology pattern in which these
two endpoint characters occur with opposite signs is
excluded by the same equality in $K_0(X)$.

For rank $r=2$ there is a further exclusion with equal
signs: there are no such $A,T$ satisfying
\[
 \operatorname{ch}(A)+\operatorname{ch}(T)
                       =3e^{-B}(1-qA_1)(1-qA_2).
 \tag{F1R15.2}
\]
Thus both parities of two-cohomology placements are covered
when the bundle endpoint has rank two. The equal-sign
rank-three case is not asserted here.

\paragraph{The actual operation clearing the divisor twist.}
Choose an integer $m>0$ such that $mB$ is the first Chern
class of a line bundle. For the actual multiplication
isogeny $p=[m]:X\to X$, the class $p^*B=m^2B$ is integral.
Choose a line bundle $M$ with $c_1(M)=p^*B$ and apply the
actual exact functor
\[
                         \mathcal F\longmapsto p^*\mathcal F\otimes M.
 \tag{F1R15.3}
\]
It preserves the ranks and local freeness of $A$ and
changes the displayed characters by the precise formulas
\[
 q'=m^4q,\qquad \ell'=m^2(\ell+B),\qquad
 e^{c_1(M)}p^*\bigl(e^{-B}v_q\bigr)=v_{q'},\qquad
 \operatorname{ch}(p^*T\otimes M)=e^{\ell'}.
 \tag{F1R15.4}
\]
Thus the argument below takes place on the actual vector
bundle $A'=p^*A\otimes M$. No formal rational line bundle
or unproved preservation of bundle rank under a change
of character is used.

\paragraph{The exact Newton identities in codimension four.}
Put $C=A_1+A_2$ and $\nu=A_1A_2$; then
$C^2=2\nu\ne0$. On the fourfold, the Newton identity
relating total Chern class to full character is
\[
 \log c(x)=\sum_{j=1}^4(-1)^{j-1}(j-1)!
                      \operatorname{ch}_j(x).
 \tag{F1R15.5}
\]
For the actual virtual class with character
$Nv_{q'}=N-Nq'C+N(q')^2\nu$, it gives
\[
 \begin{split}
 \log c(Nv_{q'})&=Nq'C-6N(q')^2\nu,\\
 c(Nv_{q'})&=1+Nq'C+N(N-6)(q')^2\nu.
 \end{split}
 \tag{F1R15.6}
\]
Here $c(Nv_{q'})$ denotes the total Chern class determined
by this character; the virtual class itself is the actual
difference or sum in (F1R15.1) or (F1R15.2) after applying
(F1R15.3). The exponential rank-one character gives
$c(T')=1+\ell'$ by the same Newton identity.

For (F1R15.1), the actual relation is
$[A']=[A'-T']+[T']$, with $N=r-1$. Multiplicativity
of total Chern class therefore gives
\[
 c(A')=\bigl(1+Nq'C+N(N-6)(q')^2\nu\bigr)(1+\ell'),
 \qquad
                c_4(A')=N(N-6)(q')^2\nu.
 \tag{F1R15.7}
\]
The codimension-three term of the first factor is zero,
so no term involving $\ell'$ enters $c_4(A')$. For
$N=1,2$ its coefficient is respectively $-5,-8$.
Both are nonzero. But an actual vector bundle of rank
two or three has $c_4=0$. This proves (F1R15.1).

For (F1R15.2), the actual relation is instead
$[A']=[A'+T']-[T']$, so
\[
 c(A')=\frac{1+3q'C-9(q')^2\nu}{1+\ell'}.
 \tag{F1R15.8}
\]
Its relevant components are
\[
 c_3(A')=-(\ell')^3-3q'C\ell',\qquad
 c_4(A')=(\ell')^4+3q'C(\ell')^2-9(q')^2\nu.
 \tag{F1R15.9}
\]
Rank two forces both to vanish. Multiplying the first
equality by $\ell'$ shows that the first two terms in
the second sum are zero, leaving
$c_4(A')=-9(q')^2\nu\ne0$, a contradiction.

\paragraph{Actual coupled-object and framework scope.}
For $H^0(E)=A$ and $H^1(E)=T$, the actual canonical
triangle gives $[E]=[A]-[T]$, so (F1R15.1) excludes
every connecting class, including nonzero ones. For
any other two cohomology degrees the canonical triangle
relations give their actual parity signs. A global shift
multiplies the whole character by its actual sign and
reduces the same-sign or opposite-sign case to the two
displayed equations.

Equations (F1R15.3)--(F1R15.9) use only actual isogeny
pullback, line tensor, the relation in the Grothendieck group,
and the universal Newton identities. These are the
$\lambda$-ring and Adams framework of F04; the codimension
scalars $p^*|_{H^{2j}}=m^{2j}$ are written explicitly.
The pullback has its actual normalized-trace retraction,
and its line tensor has the Poincar\'e biextension
conjugation already specified in R08--R09. No new
semiregularity premise is imported. The obstruction
concerns an actual low-rank bundle endpoint; it does not
assert the same vanishing of high Chern classes for an
arbitrary torsion-free sheaf or perfect complex.
